\chapter{Open Questions}

The narrow v1 contract --- one volume, one pinned state, namespace
plus logical size, fallback for everything else --- is complete.
What follows are the named investigation lanes the project has
opened deliberately, with enough scope clarity that they can be
worked on without widening v1's support claims by accident.

The chapter is organised by R2 lane rather than by topic. A lane is
a unit of work whose entire scope --- source class, semantic mode,
oracle, support claim --- can be specified before the work starts.
Lanes are isolated from each other; finishing one does not entail
finishing any other.

\section{R2-A: Physical Size Per File}

The most visible WizTree feature R1 does not yet provide.

\textbf{The question.} For each inode, what number of bytes does
APFS allocate to its content on disk?

\textbf{Why now.} Two of the candidate sources are already decoded
by the v1 body parser. \texttt{j\_dstream\_t.alloced\_size} is
emitted by the inode body decoder; \texttt{j\_file\_extent\_*}
records are walked by the FS-tree walker and counted in the family
histogram. The chapter~6 cursor rule already gives us
correct field reads. What is missing is the precedence rule that
turns the candidates into one number per inode, with an oracle
that says when the number is right.

\textbf{Candidate sources.}

\begin{description}
  \item[\texttt{j\_dstream\_t.alloced\_size}.] The dstream's claim
    about allocated bytes. Easy to read, but APFS clone records
    and snapshot-retained data can make the dstream's view
    incomplete.
  \item[\texttt{j\_file\_extent\_val\_t} records.] Per-inode
    extents, walked from the FS tree. Summing the
    \texttt{length\_and\_flags} of an inode's extents gives the
    inode's apparent on-disk extent.
  \item[Extent reference tree.] The container-level structure that
    tracks block ownership across clones and snapshots. Required
    when the question is "bytes uniquely owned by this inode" but
    not when the question is "bytes allocated to this inode's
    extents."
\end{description}

\textbf{Oracle.} \texttt{st\_blocks * 512}, for cases that have
been validated case by case. The match is not always exact ---
APFS may allocate metadata blocks that \texttt{st\_blocks} does not
count, and clone behaviour means two inodes' allocated bytes do not
sum to disk usage. The probe records both numbers and the
divergence, rather than asserting equality.

\textbf{Entry path.} SR-019 (source review of allocation
semantics: dstream vs extents vs extent reference tree) then
EX-22 (same-run fixture: ordinary, sparse, clone, hard link,
symlink, compressed; \texttt{st\_blocks * 512} per inode; raw
allocated number per source candidate; divergence breakdown). Pass
condition: for every case, exactly one source candidate matches
\texttt{st\_blocks * 512}, and the precedence rule that picks it is
stated in SR-019.

\textbf{What stays out.} Exclusive bytes (bytes only this inode
owns, after accounting for clones and snapshots) and shared bytes
(bytes shared across inodes or snapshots) are separate metrics
with separate oracles. They are not R2-A. The extent-reference
tree is decoded only as far as R2-A's allocated-size question
requires.

\section{R2-B: Snapshot-Assisted Scanning}

The mechanism by which a fast indexer can scan a live machine
honestly.

\textbf{The question.} If the indexer takes a local APFS snapshot
of a live volume, mounts it read-only, and walks it through the
existing fallback, does the result match what the same walker would
have emitted on the live volume at the snapshot moment?

\textbf{Why now.} The snapshot mechanism is the only published
macOS API that pins a coherent state of a live volume. The
fallback walker is shape-equivalent to raw mode on the proof
fixture (chapter~11). Combining the two lets the indexer scan a
live volume without depending on the raw-mode allowlist, and
without needing to handle live churn (chapter~3).

\textbf{Oracle.} The live-volume scan at the snapshot moment is its
own oracle: the snapshot is, by construction, a copy of the live
state at the moment it was taken. The probe compares the snapshot
walk to the live walk, taken simultaneously, for unchanged data.

\textbf{Entry path.} SR-020 (source review of snapshot API:
\texttt{fs\_snapshot\_create}, snapshot identity, mount lifecycle,
cleanup) then EX-23 (live-vs-snapshot shape parity fixture). Pass
condition: every \texttt{NamespaceEntry} and every
\texttt{DirectoryAggregate} that the snapshot walk produces is
present and identical in the live walk for paths that did not
change during the probe.

\textbf{What stays out.} Snapshot-retained bytes (bytes a snapshot
holds against deletion) is a separate metric and a separate
oracle. Sealed-system content (the parts of the macOS startup
volume that user snapshots cannot mount) is a separate access
class. Neither belongs in R2-B.

\section{R2-C: Scanner Ergonomics}

The lane that does not change semantics, only the surface.

\textbf{The question.} Which parts of the scanner's user-facing
surface --- backend selection, progress reporting, embeddability
in a native UI, batch syscalls --- can be improved without moving
the emission contract?

\textbf{What has landed.} Bulk syscalls for the fallback path
(\texttt{getattrlistbulk}, chapter~12) and stderr-streamed
progress events with one update per directory boundary. A native
macOS shell (\texttt{app/}) that runs the scanner as a subprocess
and hosts the result in a treemap.

\textbf{What is still in scope.} A native NSMenu for row context
operations (Reveal in Finder, Move to Trash, Copy Path); a Top-N
sidebar so a treemap with a long tail is browsable; a search
field that filters rows by stored UTF-8 name; bundling the
scanner binary into the app so the user does not need to run
\texttt{cargo build} themselves; code signing and notarisation
for a distributable build.

\textbf{What stays out.} The shape of \texttt{ParserOutput} does
not move in R2-C. A new column would be R2-A or R2-B. R2-C only
makes the existing columns easier to consume.

\section{Beyond R2}

Three areas are deliberately deferred past R2. Each is large
enough to warrant its own gate and its own oracle work; each is
unsafe to open opportunistically.

\paragraph{Encryption.} FileVault-encrypted volumes and
hardware-backed encryption have separate threat models, separate
privilege requirements, and separate validation problems. The
parser currently fails closed on encryption flags; opening
encryption support requires deciding which classes of key access
the product is willing to use and which it is not.

\paragraph{Live boot disk.} The startup disk on modern macOS is a
sealed system volume joined to a data volume by firmlinks, with a
snapshot pinning the system content. Raw scanning of the System
volume is constrained by the seal; raw scanning of the Data volume
is constrained by the live-churn problem; the merged user-visible
view depends on firmlink table semantics that are themselves
not fully public. Opening live boot support means building each
of those subsystems' oracle.

\paragraph{Boot-root merged namespace.} The Finder-visible
\texttt{/} on a modern Mac is not any single APFS volume's
namespace. It is the System volume's content joined to the Data
volume's content through firmlinks, plus the current snapshot
overlay. The product mode for this is distinct from
raw-single-volume; it cannot share the same oracle, and its
support claims must be stated separately.

\paragraph{Persistent incremental cache.} The work of chapter~10
turns into product only after the full set of cache-invalidation
conditions has been enumerated, a crash-safe cache storage design
exists, and an incremental scan has been validated against a
fresh full reparse. The R2 lanes deliberately precede this work;
a cache built on top of incomplete metrics or unstable support
matrices would propagate its instability.

\section{The Manual's Own Future}

The chapter ends where the manual ends. Each lane above will, when
it lands, become its own chapter: a thesis, the evidence, the
precedence rule or the oracle pairing, the per-case table, the
limits of what was checked. The chapter list will lengthen; the
voice should not change.

The discipline that produced R1 produced something narrow but
defensible. The discipline that produces R2 will produce something
broader and still defensible. Anything written here that survives
the next year of work survives because it was written carefully.
Anything that does not survive was an error worth recording so the
next attempt does not repeat it.
